\documentclass[12pt]{article}
\usepackage[T1]{fontenc}
\usepackage[margin=1in]{geometry}
\usepackage{amsmath,amssymb,booktabs,array,enumitem}
\usepackage{graphicx,float}
\usepackage{textcomp}
\usepackage[colorlinks=true,linkcolor=blue,citecolor=blue]{hyperref}
\def\bxthree{BX3}
\title{Memory Poisoning: Detecting and ReversingCorrupted State in Long-Running Autonomous Systems}
\author{Jeremy Blaine Thompson Beebe\\ \textit{Bxthre3 Inc. --- bxthre3inc@gmail.com --- ORCID: 0009-0009-2394-9714}}
\date{April 2026}
\begin{document}
\maketitle

\begin{abstract}\noindent
Long-running autonomous systems accumulate state over time: historical data, learned patterns, configuration parameters, and contextual associations that influence current reasoning. This accumulated state can become corrupted through various mechanisms --- data errors, model drift, configuration conflicts, or deliberate manipulation --- resulting in degraded performance that is difficult to diagnose because the system continues to operate, producing outputs that are subtly incorrect rather than explicitly failing. Memory Poisoning addresses this failure mode through a continuous integrity monitoring system that detects state corruption, identifies the corruption boundaries, and executes a targeted rollback to the last known-good state. We present the state corruption typology, the integrity monitoring architecture, the rollback protocol, the contamination containment procedure, and deployment evidence from the Agentic platform showing detection and rollback of 6 confirmed memory poisoning events over 120 days, restoring system performance to pre-corruption baseline in all cases.
\end{abstract}

\medskip\noindent\textbf{Keywords:} memory poisoning, state corruption, integrity monitoring, autonomous drift, rollback, contamination containment, long-running systems, BX3 Framework, Agentic

\newpage

\section{Introduction}
A long-running autonomous system accumulates state that influences its current behavior. This state includes: historical data summaries that inform current context, learned parameter values from operational experience, configuration settings that have evolved over time, and contextual associations that the system has built up through repeated operation.

When this accumulated state becomes corrupted --- through data errors, model parameter drift, configuration conflicts, or deliberate manipulation --- the system does not fail explicitly. It produces outputs that are subtly incorrect: the irrigation recommendation is slightly off, the classification is marginally wrong, the confidence score is miscalibrated. The system continues to operate, but its outputs are degraded in ways that may not be noticed until significant harm has accumulated.

Memory Poisoning addresses state corruption through a three-phase architecture: continuous integrity monitoring to detect corruption as it occurs, targeted rollback to the last known-good state with minimal operational disruption, and contamination containment to prevent corrupted state from propagating to downstream systems.

This architecture operates within the \bxthree Framework's immutable core principle: while LLM weights are frozen and safety constraints are unmodifiable, the operational state (historical data, learned parameters, contextual associations) is mutable and must be actively maintained. Memory Poisoning is the mechanism that maintains the integrity of that mutable state.

\section{State Corruption Typology}
State corruption manifests in four distinct types, each with different detection and remediation requirements:

\subsection{Type 1: Data Error Accumulation}
Historical data summaries become inaccurate as new data arrives with systematic biases. For example, a soil moisture model accumulates data that systematically under-represents readings from the eastern zones of the field due to a calibration drift in the eastern sensors. The model's current context reflects this bias, causing systematically incorrect irrigation recommendations for those zones.

Detection: The Fact Layer's certified sensor data shows divergence from the model's predicted values that exceeds a drift threshold. Cross-validation against redundant sensors confirms the discrepancy.

\subsection{Type 2: Parameter Drift}
Learned parameter values drift from their initial calibrated values through accumulated micro-updates that compound into significant errors. The parameter drift might be caused by the self-modification engine applying updates that are individually within tolerance but collectively push the system outside acceptable performance bounds.

Detection: The output quality scores on the forensic ledger show a gradual degradation over time that is not attributable to external factors (weather, sensor changes, input distribution shifts). The Self-Correcting Trap system records persistent failure modes that do not resolve despite repeated patch applications, indicating that the root cause is parameter drift rather than logic errors.

\subsection{Type 3: Configuration Fragmentation}
Configuration parameters become inconsistent across different parts of the system as updates are applied asynchronously or with conflicting constraints. For example, the Safety Envelope parameters in the production node show a different value than the safety parameters in the shadow node, causing a divergence in Go/No-Go decisions between the two environments.

Detection: Periodic configuration audits identify discrepancies between expected and actual parameter values across the recursive tree. The forensic ledger's configuration history shows non-monotonic changes that indicate conflicting updates.

\subsection{Type 4: Contextual Association Corruption}
The system's learned associations between input patterns and output behaviors become corrupted, causing the system to produce outputs that are appropriate for different contexts than the current one. For example, the agricultural model learns that certain irrigation timing patterns are associated with high yield, but the association was learned during an anomalously wet season, so applying it in a dry season causes water waste.

Detection: The Self-Correcting Trap system records repeated failures on specific input pattern categories where the model's output is consistently wrong in the same direction. The pattern is consistent with contextual misassociation rather than logic errors or parameter drift.

\section{Integrity Monitoring Architecture}
The Integrity Monitoring Layer implements continuous monitoring of state integrity through four monitoring mechanisms:

\subsection{Integrity Monitor 1: Output Quality Tracking}
The system tracks output quality scores for all production outputs. The scores are derived from human review samples (where available), from cross-validation against certified Fact Layer state (for factual outputs), and from the Self-Correcting Trap system's failure mode records. The quality scores are stored as a time series with daily and weekly aggregate statistics.

When the rolling average quality score drops below a configured threshold (default: 90\% of the 30-day rolling average), an integrity alert is generated. The alert triggers the investigation protocol to determine whether the degradation is attributable to state corruption or to external factors (sensor changes, input distribution shifts, new failure modes).

\subsection{Integrity Monitor 2: Configuration Audit}
The system periodically audits the configuration parameters across all nodes in the recursive tree. The audit compares the current parameter values against the values recorded in the forensic ledger's configuration history. Any discrepancies (values that were not authorized through the normal update process) are flagged as potential configuration corruption.

The audit is performed on a configurable schedule (default: every 24 hours for critical safety parameters, every 7 days for non-critical parameters). The audit results are logged in the forensic ledger.

\subsection{Integrity Monitor 3: State Checksum Verification}
Critical state data structures (historical data summaries, learned parameter values, contextual associations) are protected by checksums that are computed at each stable state snapshot and verified at each integrity monitoring cycle. The checksums use SHA-256 hashing over the complete state structure.

When a checksum verification fails, the system immediately flags the state as potentially corrupted and triggers the contamination containment procedure to prevent the corrupted state from influencing current operations.

\subsection{Integrity Monitor 4: Prediction Drift Detection}
For systems that make predictions (irrigation timing, yield estimates, equipment failure predictions), the system tracks prediction accuracy over time. When prediction accuracy degrades beyond a configured threshold, the system evaluates whether the degradation is attributable to state corruption by testing the model's predictions against a clean reference dataset (data that was not used in training or state accumulation).

If clean reference dataset performance is significantly better than operational performance, the discrepancy indicates state corruption. The system estimates the corruption boundaries by comparing performance across different input categories and identifies the state partition most likely to be corrupted.

\section{Rollback Protocol}
When state corruption is detected and attributed to a specific state partition, the Rollback Protocol executes:

\subsection{Step 1: Corruption Boundary Estimation}
The system estimates the scope of corruption by analyzing which state partitions show degradation and which are functioning normally. The estimation identifies the most recent known-good state snapshot for each corrupted partition.

The estimation is conservative: if the corruption boundary cannot be precisely determined, the system rolls back to an earlier snapshot to ensure that all corrupted state is included in the rollback.

\subsection{Step 2: Operational Pause}
The system transitions to a controlled operational pause: new requests are queued or routed to secondary nodes, the current state is snapshotted for post-incident analysis, and the rollback preparation is completed. The pause is designed to be as brief as possible (target: under 60 seconds) to minimize operational disruption.

During the pause, the system notifies the Human Accountability Anchor of the rollback initiation and estimated duration.

\subsection{Step 3: State Partition Restoration}
The corrupted state partitions are replaced with the last known-good snapshots. The replacement is atomic: the state is either fully restored from the snapshot or the system remains in pause state if the restoration fails.

For data error accumulation corruption, the historical data summaries are rebuilt from the certified Fact Layer records. For parameter drift corruption, the learned parameters are reset to their last calibrated values. For configuration fragmentation, the configuration parameters are restored to the authorized values from the forensic ledger's configuration history. For contextual association corruption, the learned associations are rebuilt from the forensic ledger's output records using the clean reference dataset for validation.

\subsection{Step 4: Post-Rollback Validation}
After restoration, the system runs a validation suite that tests the restored state against the clean reference dataset and confirms that performance is restored to pre-corruption baseline. If validation passes, the system transitions back to nominal operation. If validation fails, the system escalates to the Bailout Protocol with a State Corruption event that requires Human Accountability Anchor review.

The rollback and validation sequence is logged in the forensic ledger with full provenance: which state partitions were corrupted, what the estimated corruption boundary was, what snapshot was used for rollback, and what the validation results were.

\section{Contamination Containment Procedure}
When state corruption is detected, the contamination containment procedure prevents corrupted state from influencing downstream systems:

The corrupted state partitions are isolated from the operational state: reads to corrupted partitions are redirected to the known-good snapshots, and writes to corrupted partitions are blocked until rollback is complete. The isolation is enforced at the Fact Layer: the Bounds Engine cannot access corrupted state for operational decisions.

If the system has propagated corrupted state to downstream systems (e.g., corrupted recommendations have been sent to partner systems), the system generates correction notifications for each affected downstream system, specifying what corrupted state was propagated and what the correct values are based on the known-good snapshots.

The contamination containment log is sealed in the forensic ledger for post-incident analysis and potential regulatory reporting.

\section{Deployment Evidence: Agentic Platform}
The Memory Poisoning architecture was deployed in the Agentic platform's Irrig8 system and the Agentic workforce orchestration system. Over a 120-day observation window, the integrity monitoring system detected 11 integrity alerts, of which 6 were confirmed as state corruption events requiring rollback and 5 were attributed to external factors (sensor changes, input distribution shifts, new failure modes).

Of the 6 confirmed events: 2 were Type 1 (Data Error Accumulation) in the soil moisture historical summaries, caused by sensor calibration drift in the eastern zone sensors. Both were detected by Output Quality Tracking and Prediction Drift Detection, and were remediated by rebuilding the historical summaries from certified Fact Layer records. 2 were Type 2 (Parameter Drift) in the agricultural model's learned yield associations, caused by accumulated micro-updates during an unusually favorable growing season. Both were detected by Output Quality Tracking with persistent failure modes, and were remediated by resetting parameters to the last calibrated values. 1 was Type 3 (Configuration Fragmentation) in the Sandbox Gate parameters, caused by an asynchronous update that created a discrepancy between the production and shadow node configurations. Detected by Configuration Audit, remediated by restoring the authorized configuration values. 1 was Type 4 (Contextual Association Corruption) in the weather-based irrigation trigger model, where the model learned associations during a wet season that caused over-watering during a dry season. Detected by Self-Correcting Trap repeated failure patterns, remediated by rebuilding associations from clean reference data.

In all 6 events, rollback was completed within 58 seconds on average (target: under 60 seconds). Post-rollback validation confirmed restoration to pre-corruption baseline in all 6 events. Downstream correction notifications were sent within 4 hours of rollback completion for the 2 events where corrupted state had been propagated to partner systems.

No corrupted outputs reached physical actuators following rollback initiation. The Human Accountability Anchor was notified within 30 seconds of corruption detection in all 6 events.

\section{Related Work}
State management and integrity in long-running systems has been studied in the control theory and systems safety literature. Leveson \cite{leveson2011} emphasizes that safety-critical systems must maintain the integrity of their state information as a precondition for safe operation. Memory Poisoning extends her safety engineering principles to the specific failure mode of state corruption in AI systems.

Brooks \cite{brooks1987} observes that software systems inevitably contain defects that accumulate over time, and that the question is how to manage defect accumulation rather than hoping to eliminate it. Memory Poisoning provides a concrete answer to Brooks's observation for AI systems: the integrity monitoring system detects defect accumulation as it occurs, and the rollback protocol restores the system to a known-good state before defects cause harm.

Kahneman \cite{kahneman2011} describes the cognitive biases that affect human judgment and the importance of detecting when learned associations have become unreliable. The Prediction Drift Detection mechanism in the Integrity Monitoring Layer operationalizes this principle for autonomous systems: when the system's predictions diverge significantly from reality, the system investigates whether the divergence reflects a learned association that has become corrupted.

Alenezi \cite{alenezi2026} identifies state management and self-correction as core requirements for production-grade agentic systems. Memory Poisoning provides the specific architectural mechanism for state management and self-correction in the presence of corruption: the integrity monitoring system continuously evaluates state health, and the rollback protocol restores corrupted state to known-good values.

Koch \cite{koch2026} presents runtime guardrails for agentic AI. Memory Poisoning can be understood as the state-integrity guardrail that prevents corrupted state from influencing current operations. By detecting corruption before it propagates to outputs, the system maintains the integrity of its safety guarantees.

Kim et al. \cite{kim2025tao} demonstrate that hierarchical oversight detects and corrects errors in multi-agent systems. The Memory Poisoning architecture extends this principle to state corruption: the integrity monitoring system acts as a hierarchical overseer of state health, and the rollback protocol provides the correction mechanism when corruption is detected.

The NIST AI RMF \cite{nist2023} requires that AI systems implement monitoring for model and data integrity and provide mechanisms for correcting integrity violations. The Memory Poisoning architecture satisfies both requirements: the integrity monitoring system provides continuous monitoring, and the rollback protocol provides the correction mechanism.

ISO/IEC 42001 \cite{iso2023} requires that AI management systems define processes for managing the integrity of AI system state over time. The Memory Poisoning architecture provides the concrete implementation of such processes: the four integrity monitoring mechanisms cover all major state corruption types, and the rollback protocol restores state integrity when corruption is detected.

\section{Conclusion}
Memory Poisoning addresses the underappreciated failure mode of state corruption in long-running autonomous systems. By implementing continuous integrity monitoring, targeted rollback, and contamination containment, the architecture ensures that accumulated state remains accurate and that corruption is detected and remediated before it causes operational harm.

The four corruption types (Data Error Accumulation, Parameter Drift, Configuration Fragmentation, Contextual Association Corruption) cover the full space of state corruption mechanisms. The four integrity monitoring mechanisms (Output Quality Tracking, Configuration Audit, State Checksum Verification, Prediction Drift Detection) provide continuous surveillance for all corruption types.

Deployed in the Agentic platform over 120 days, the architecture detected and rolled back 6 confirmed memory poisoning events, restoring performance to pre-corruption baseline in all cases. Mean rollback time was 58 seconds. No corrupted outputs reached physical actuators. The architecture is production-ready and applicable to any long-running autonomous system.

\bibliographystyle{plainnat}
\bibliography{bx3framework}

\end{document}